\begin{titlepage}

    \newcommand{\HRule}{\rule{\linewidth}{0.5mm}} 

    \begin{center}
    
        \textsc{\LARGE Old Dominion University}\\[1.0cm]
        \textsc{\Large CS 895: LLM Architecture \& Applications}\\[0.5cm] 
        \HRule \\[0.1cm]
        { \huge \bfseries Prompting Strategies and Their Effect on Model Behavior}\\[0.3cm] 
        { \Large \bfseries Homework 2}\\[0.3cm] 
        \HRule \\[0.1cm]

      
        \begin{minipage}{0.4\textwidth}
            \begin{flushleft} \large
                \emph{By:} \\
                Nahom M. Birhan \\ 
                ID: nbirh002
            \end{flushleft}
        \end{minipage}
        ~
        \begin{minipage}{0.4\textwidth}
            \begin{flushright} \large
                \emph{Lecturer:} \\
                Prof. Nazzal
            \end{flushright}
        \end{minipage}\\[0.5cm]

        \emph{Date Due:} March 4, 2026, \emph{Date Submitted:} \today \\[1.0cm] 

    \end{center}
   \vspace{0.2cm}
    \begin{abstract}
        \noindent This report investigates prompting strategies and their effect on large language model behavior across three axes. Prompts, examples, and evaluation sets were constructed with LLM assistance. The first examines in-context learning for sentiment classification on sarcastic product reviews, comparing zero-shot, one-shot, and few-shot prompting with Gemini 2.5 Flash and Claude Haiku; in-context examples improved accuracy from 0.80 to 0.90 but few-shot did not uniformly outperform one-shot across models. The second compares direct-answer and chain-of-thought prompting on twelve arithmetic word problems, finding that CoT reduced Gemini's accuracy from 0.83 to 0.75 through propagated arithmetic errors while slightly improving Claude's precision and F1 at equal accuracy. The third measures prompt sensitivity through eight controlled perturbations of the few-shot baseline, including example reordering, instruction rephrasing, constraint addition and removal, reduced example count, contradictory labels, role changes, and label-space expansion; Gemini held at 0.90 for four variants and dropped to 0.80 for the remaining four, while Claude was more volatile, falling to 0.70 on constraint removal but rising to 0.90 with fewer examples, the only configuration where a perturbation improved over its 0.80 baseline. These findings demonstrate that model behavior is highly prompt-dependent and that strong performance on a fixed prompt does not guarantee robustness, underscoring the need for systematic prompt evaluation in deployable systems. Source code: \texttt{https://github.com/NahomMA/llm-prompt-sensitivity-and-cot}.
    \end{abstract}

\end{titlepage}